\documentclass{article}
\usepackage[utf8]{inputenc}
\usepackage[T1]{fontenc}
\usepackage{geometry}
\usepackage{graphicx}
\usepackage{float}
\usepackage{hyperref}
\usepackage[english,indonesian]{babel}

\title{Optimasi Pemrosesan Template LaTeX Menggunakan Pendekatan Berbasis AI}
\author{Academic Support Agent}
\date{\today}

\begin{document}
\maketitle

\begin{abstract}

\end{abstract}

\chapter{PENDAHULUAN}


\chapter{TINJAUAN PUSTAKA}


% Original content (commented out):
% % Template LaTeX para o III CBCFD
% Versao para o artigo
\documentclass[a4paper,12pt]{article}
% Configurações para o template - Não excluir!
% Pacotes LaTeX necessários:
% babel, inputenc, geometry, fontenc, indentfirst, natbib,
% amsmath, amsfonts, amssymb, graphics, graphicx, epsfig, color
% hyperref, textcomp, fancyhdr, fancyvrb, caption
\input{./config/cbcfd_config}

% INICIO DO DOCUMENTO LaTeX
\begin{document}

% Coloque o título aqui
\title{O título do trabalho do III Congresso de Fluidodinâmica Computacional 
deve ser escrito com o comando $\backslash$\MakeLowercase{title}$\{\}$}

% Coloque os autores aqui
\author{A. B. C. Dario$^{1}$, E. F. da Guia$^{2}$ e G. H. Informa Jr$^{3}$}

% Coloque as filiações
\address{
$^{1}$ Universidade Estadual de Campinas,
       Faculdade de Engenharia Química \\
$^{2}$ Universidade Federal de Campina Grande,
       Departamento de Engenharia de Materiais \\
$^{3}$ Universidade Federal do Rio de Janeiro,
       Escola de Química \\
\email{dario@wxyz.br}
}

% Coloque aqui o resumo
\resumo{O texto em letras minúsculas, continuando na mesma linha em um
único parágrafo com letra Times New Roman 12 e espaçamento simples. A
formatação automática do resumo está definida no comando
$\backslash$resumo$\{\}$ em \LaTeX. O texto no resumo não deverá conter
referências bibliográficas, sendo limitado a 800 caracteres contando os
espaços. Este \textit{template} foi escrito em \LaTeX e deve-se usar os
comandos descritos no arquivo fonte \texttt{.tex}.
}

% Início das seções
\section{Introdução}
O título da seção é um cabeçalho de primeira ordem, deverá ser numerado
com algarismos arábicos, alinhados junto à margem esquerda, com letras
maiúsculas e em negrito. Deverá ser separado por um espaço duplo antes
e depois do texto. A formatação está definida pelo comando
$\backslash$section$\{\}$.

O texto deverá ser digitado em espaço simples para todo o trabalho e
espaço duplo entre parágrafos. A página deverá ser configurada no
formato A4 ($210mm$ x $297mm$), com orientação ``retrato'', em uma
coluna, com margem superior $40~mm$ e margens inferior, direita e
esquerda de $25~mm$. A fonte deverá ser a Times New Roman, tamanho 16
no título, 14 nos cabeçalhos de primeira e segunda ordem, 12 no texto
e cabeçalhos de terceira ordem e 12 nas notas de referência dos autores
abaixo do título. Recomenda-se que o trabalho completo seja escrito
neste mesmo arquivo fornecido pela comissão organizadora.

Os trabalhos redigidos em português, inglês ou espanhol deverão ter a
mesma formatação. A primeira linha de cada parágrafo deverá iniciar com
uma tabulação de $1~cm$ da margem esquerda, com alinhamento justificado.
O trabalho completo deverá ter no máximo de 04 (quatro) páginas.

\subsection{Segundo Cabeçalho}
Os cabeçalhos de segunda ordem deverão ser alinhados junto à margem
esquerda, digitados em negrito, sendo maiúscula apenas a primeira letra
de cada palavra com separação de um espaço do texto anterior. A formatação
está definida pelo comando $\backslash$subsection$\{\}$.

\subsubsection{Terceiro cabeçalho}
O cabeçalho de terceira ordem inicia após tabulação de $1~cm$ da
margem esquerda, sem numeração, com caracteres sublinhados, sendo
maiúscula apenas a primeira letra da primeira palavra. O texto deverá
começar na mesma linha do cabeçalho. A formatação está definida pelo
comando $\backslash$subsubsection$\{\}$.

\section{Figuras}
As figuras poderão ser coloridas e deverão ser inseridas no corpo do
trabalho, tão próximas quanto possível das citações, no texto, sobre
elas. As figuras deverão ser centralizadas, sem exceder o tamanho
limitado pelas margens da página.

Cada figura deverá ter um título numerado em algarismos arábicos. Os
títulos deverão ser centralizados na parte inferior das mesmas,
separados por espaço simples antes e duplo após o texto e digitados
como: \figref{fig:alg} -- Título da figura com ponto final. No texto,
elas deverão ser mencionadas, por exemplo, da seguinte forma: ``conforme
mostra a \figref{fig:alg}...''.  Estes comandos estão
definidos de acordo com a nomenclatura de \LaTeX, usando 
$\backslash$figref$\{\}$.
\begin{figure}[htb]
\centering
\includegraphics[width=\textwidth]{./images/FiguraExemplo}
\caption{Vetores velocidade coloridos com a temperatura.}
\label{fig:alg}
\end{figure}

\section{Tabelas}
As tabelas deverão ser centralizadas e posicionadas após e próximas às
suas respectivas citações no texto, com os títulos centralizados na
parte superior das mesmas com espaço duplo antes e simples após o texto
e digitados como: \tabref{tab:tabela} -- Título da tabela sem ponto
final. As tabelas deverão ser mencionadas no texto, por exemplo, da
seguinte forma ``...o resultado do planejamento pode ser encontrado na
\tabref{tab:tabela}''. Estes comandos estão
definidos de acordo com a nomenclatura de \LaTeX, usando 
$\backslash$tabref$\{\}$. As unidades de medida correspondentes a todos
os termos deverão ser claramente indicadas, preferencialmente no S.I.
\begin{table}[!htb]
  \caption{Serviços Térmicos}
  \begin{center}
    \begin{tabular}{| l p{4cm} p{4cm} p{4cm} |}
    \hline
    Trocador & Vazão da Corrente Quente ($kg/s$)
    &  Temperatura da corrente quente ($^{o}C$)
    & Vazão de água de resfriamento($kg/s$)\\ \hline
    E-101 & 12,0 & 80  & 11,8 \\
    E-102 & 12,0 & 120 & 28,7 \\
    \hline
    \end{tabular}
  \end{center}
  \label{tab:tabela}
\end{table}

\section{Equações}
As equações deverão ser escritas em itálico, após uma tabulação (1 cm),
com numeração consecutiva entre parênteses, rente à margem direita e
espaço duplo antes e após a equação. Deverão ser mencionadas no texto,
por exemplo, da seguinte forma: “... a substituição da
\eqnref{eq:radiacao} na Equação 3 (referenciada em \LaTeX) fornece...”.
Equações com mais de uma linha deverão ser numeradas na última linha,
entre parênteses e rente à margem direita. Estes comandos estão
definidos de acordo com a nomenclatura de \LaTeX, usando 
$\backslash$eqnref$\{\}$.
\begin{equation}
\varepsilon T_{c,i} + (1-\varepsilon) T_{h,i} - T_{h,o} = 0
\label{eq:radiacao}
\end{equation}

\section{Nomenclatura}
Caso necessário, a lista de notações e símbolos utilizados, assim como
suas unidades de medida, deverão ser relacionados antes das referências
bibliográficas por ordem alfabética.

\section{Referências}
As referências bibliográficas deverão ser citadas no texto, através do
último sobrenome do autor e do ano de publicação, o qual deverá estar
entre parênteses, conforme os exemplos a seguir: ``Os trabalhos de Reh
et al. (1977) e Souza (1996) mostraram...'', ou ``... tem sido mostrado
(Reh et al., 1977; Souza, 1996)''. No caso de dois autores, ambos
deverão ser citados, exemplo: ``... segundo Rhodes e Geldart (1995)...''
Em caso de três ou mais autores, deverá ser citado o sobrenome 
do primeiro autor seguido da expressão ``et al.''. Estes comandos estão
definidos de acordo com a nomenclatura de \LaTeX.

Trabalhos publicados no mesmo ano e pelos mesmos autores deverão usar
as letras a, b, c, junto ao ano, exemplo: ``Os trabalhos de Campos et
al. (1996 a, b) mostraram...''.

A lista de referências deverá incluir somente os trabalhos citados no
texto, relacionados em ordem alfabética, de acordo com o sobrenome do
primeiro autor, iniciando na margem à esquerda, com tabulação de 1 cm
a partir da segunda linha e separadas por espaço simples. As
referências deverão conter todos os sobrenomes e iniciais dos autores,
separados por ponto e vírgula. Os títulos dos periódicos deverão
aparecer abreviados e em itálico, conforme o exemplo aplicado a
\cite{Souza}, \citep{Rehetal77}, \citep{Silvaetal08a,DeitelDeitel}, 
\cite{Damian07} e \citep{Ribeiro05}:

\bibliography{bibTemplate}

\section{Disposições Gerais}
\begin{enumerate}
 \item[a.] os autores aceitam que o CBCFD 2024 tenha plenos direitos
           sobre os trabalhos enviados, podendo incluí-los nos anais,
           imprimi-los e divulgá-los, sem o pagamento de qualquer
           remuneração;
           
 \item[b.] os trabalhos serão avaliados por revisores convidados pelo
           Comitê Científico do Evento. Somente os trabalhos aceitos
           poderão ser apresentados e publicados nos anais do Congresso;

 \item[c.] o trabalho completo deverá ser enviado obrigatoriamente em
           formato PDF, segundo as datas e normas que constam na página
           oficial do Congresso:
           \url{http://2024.cbcfd.com.br/}
\end{enumerate}

\end{document}


\end{document}